\begin{enumerate}[label=\textbf{\arabic*)}, leftmargin=*]

    \item À medida que o cursor se move, a razão do divisor de tensão muda e, consequentemente, altera-se a tensão de saída.

    \item São dispositivos cuja resistência varia de acordo com a sua deformação mecânica.

    \item Os extensômetros diferem entre si quanto à direção de aplicação da força:
        \begin{itemize}
            \item \textbf{Uniaxial:} mede em apenas uma direção;
            \item \textbf{Biaxial:} mede em duas direções;
            \item \textbf{Roseta:} mede em três direções (geralmente $0^\circ$, $45^\circ$ e $90^\circ$).
        \end{itemize}

    \item Para um extensômetro metálico $\Delta R = 0{,}24\,\Omega$, enquanto para um extensômetro semicondutor $\Delta R = 12\,\Omega$. Isso significa que o semicondutor apresenta uma variação de resistência 50 vezes maior que o metálico.

    \item No NTC, à medida que a temperatura aumenta, a resistência diminui. No PTC, conforme a temperatura aumenta, a resistência também aumenta.

    \item O sensor sofre interferência no campo magnético gerado pelo seu indutor devido à aproximação ou presença de materiais metálicos.

    \item Definições de distância sensora:
        \begin{itemize}
            \item \textbf{Distância sensora nominal:} distância máxima teórica na qual o sensor pode operar, considerando um alvo padrão;
            \item \textbf{Distância sensora operacional:} distância na qual o sensor opera com segurança na prática, levando em conta as tolerâncias de fabricação e condições industriais.
        \end{itemize}

    \item O sensor capacitivo detecta materiais metálicos e não metálicos por meio da alteração no campo elétrico, enquanto o sensor indutivo detecta apenas materiais metálicos via campo magnético.

    \item Circuitos de saída PNP e NPN:
        \begin{itemize}
            \item \textbf{PNP:} a saída utiliza um transistor que chaveia o VCC (terminal positivo da alimentação), fornecendo sinal positivo para a carga;
            \item \textbf{NPN:} a saída utiliza um transistor que chaveia o GND (terminal negativo da alimentação), conectando o sinal negativo à carga.
        \end{itemize}

    \item Em um sensor NPN com saída normalmente aberta (NA), dois fios são destinados à alimentação e o terceiro controla o chaveamento do sinal negativo para a carga quando acionado. No sensor PNP normalmente fechado (NF), o acionamento interrompe o sinal positivo mantido previamente na saída.

    \item Na ligação em série, a carga só é acionada se todos os sensores forem ativados simultaneamente. Na ligação em paralelo, a carga é acionada caso qualquer um dos sensores seja ativado.

    \item A frequência máxima de operação é o número máximo de comutações (chaveamentos) por segundo que o sensor consegue realizar.

    \item Trata-se da histerese: a diferença entre a distância na qual o sensor é ativado durante a aproximação do objeto e a distância na qual ele é desativado durante o seu afastamento.

    \item O diâmetro do corpo de um sensor de proximidade (como os indutivos ou capacitivos) determina diretamente o alcance nominal de detecção ($S_n$) e o tamanho mínimo do objeto a ser detectado:
        \begin{itemize}
            \item \textbf{Alcance de detecção:} quanto maior o diâmetro da bobina ou núcleo, maior é a área da superfície ativa e a intensidade do campo gerado. Consequentemente, corpos maiores (ex.: M30) alcançam distâncias de comutação muito superiores aos de corpos menores (ex.: M8 ou M12);
            \item \textbf{Área do alvo:} para atingir a distância nominal especificada, o objeto deve possuir dimensões iguais ou superiores à superfície ativa do sensor (geralmente uma placa de aço quadrada com lado equivalente ao diâmetro do sensor). Objetos menores reduzem o alcance efetivo.
        \end{itemize}

    \item A diferença central entre sensores faceados e não faceados reside na presença de blindagem metálica ao redor da cabeça sensora, alterando o campo e a forma de instalação:
        \begin{itemize}
            \item \textbf{Sensores Faceados (\textit{Shielded} / Blindados):}
                \begin{itemize}
                    \item \textit{Construção:} possuem um anel metálico protetor envolvendo a superfície ativa até a borda frontal;
                    \item \textit{Campo elétrico/magnético:} é concentrado e direcionado apenas frontalmente (direção axial);
                    \item \textit{Instalação:} podem ser montados embutidos (rentes) em superfícies metálicas sem sofrer interferência lateral do material;
                    \item \textit{Alcance:} apresentam distância de detecção menor comparada a um sensor não faceado de mesmo diâmetro.
                \end{itemize}
            \item \textbf{Sensores Não Faceados (\textit{Unshielded} / Não Blindados):}
                \begin{itemize}
                    \item \textit{Construção:} não possuem blindagem lateral na ponta, deixando a superfície ativa sobressaída do corpo metálico;
                    \item \textit{Campo elétrico/magnético:} dispersa-se tanto frontal quanto lateralmente;
                    \item \textit{Instalação:} exigem uma área livre de metal ao redor da cabeça para evitar acionamentos falsos ou degradação do sinal;
                    \item \textit{Alcance:} oferecem maior distância de detecção (até 50\% a mais) para o mesmo diâmetro de corpo.
                \end{itemize}
        \end{itemize}

\end{enumerate}